% ADDED figure (2026-11-14): Theorem §C6c.11.2, the n! assignments of n labelled lines to the n field factors of one term are the n! terms of the product rule (n = 3); the user's PHY 513 notes, Appendix B, Derivation "The derivative counts the contractions". Generated by qft-C/tools/rulesL/figs/make_figs.sh
\tikzset{
  fdl/.style={line width=0.8pt, black},
  fdx/.style={circle, fill=black, inner sep=0pt, minimum size=3.6pt},
  fdv/.style={circle, fill=figB, draw=figB, inner sep=0pt, minimum size=5pt},
  fdlab/.style={font=\scriptsize, inner sep=2.6pt},
  fdmom/.style={->, >=Stealth, figO, line width=0.6pt},
}
\tikzset{
  fda/.style={fdl, postaction={decorate}, decoration={markings, mark=at position 0.55 with {\arrow{Stealth[length=5pt]}}}},
  fdS/.style={font=\small},
}
\tikzset{
  fline/.style={line width=0.8pt, black, postaction={decorate}, decoration={markings, mark=at position 0.56 with {\arrow{Stealth[length=5.5pt, width=4.5pt]}}}},
  sline/.style={line width=0.8pt, black, dash pattern=on 3.5pt off 2.2pt},
  blob/.style={draw, line width=0.8pt, fill=figB!12},
}
\tikzset{
  pline/.style={line width=0.8pt, black, decorate, decoration={snake, amplitude=1.6pt, segment length=5.2pt}},
}
\begin{tikzpicture}[font=\small]
  \tikzset{fac/.style={draw, line width=0.7pt, rounded corners=2pt, minimum width=10mm, minimum height=6.5mm, fill=figB!10, inner sep=1pt},
           ln/.style={draw, line width=0.7pt, circle, minimum size=6.5mm, inner sep=0pt, fill=white},
           asg/.style={line width=0.9pt, figO, -{Stealth[length=4.5pt]}},
           gh/.style={line width=0.5pt, black!35, -{Stealth[length=3.5pt]}}}
  % (a) one assignment
  \begin{scope}
    \node[fdS, anchor=west] at (-3.9,2.55) {(a) one complete set of contractions $=$ one assignment $\sigma$};
    \node[fdlab, anchor=east] at (-2.45,1.6) {lines (types $a_k$):};
    \node[fdlab, anchor=east] at (-2.45,0) {one term:}; \node[fdlab, anchor=east] at (-2.45,-0.42) {(of $\mathcal L_{\rm int}$)};
    \node[ln] (l1) at (0,1.6) {$1$}; \node[ln] (l2) at (1.6,1.6) {$2$}; \node[ln] (l3) at (3.2,1.6) {$3$};
    \node[fdlab, above] at (l1.north) {$a_1$}; \node[fdlab, above] at (l2.north) {$a_2$}; \node[fdlab, above] at (l3.north) {$a_3$};
    \node[fac] (f1) at (0,0) {$\Phi_{b_1}$}; \node[fac] (f2) at (1.6,0) {$\Phi_{b_2}$}; \node[fac] (f3) at (3.2,0) {$\Phi_{b_3}$};
    \node[anchor=east] at (f1.west) {$-\frac{1}{3!}\,c_{b_1b_2b_3}$};
    \draw[asg] (l1) -- (f2); \draw[asg] (l2) -- (f3); \draw[asg] (l3) -- (f1);
    \node[fdlab, anchor=west, align=left] at (4.1,1.15) {$\sigma = (2\,3\,1)$:\\ $b_2 = a_1,\ b_3 = a_2,\ b_1 = a_3$};
    \node[fdlab, anchor=west, align=left] at (4.1,0.25) {contributes $c_{a_3a_1a_2} = c_{a_1a_2a_3}$};
  \end{scope}
  % (b) all six
  \begin{scope}[yshift=-3.3cm]
    \node[fdS, anchor=west] at (-3.9,1.9) {(b) the $3! = 6$ assignments $=$ the $6$ terms of $\partial_{a_1}\partial_{a_2}\partial_{a_3}$ by the product rule};
    \foreach \k/\a/\b/\c in {0/1/2/3, 1/1/3/2, 2/2/1/3, 3/2/3/1, 4/3/1/2, 5/3/2/1}{
      \begin{scope}[xshift=\k*2.05cm - 1.6cm]
        \foreach \j in {1,2,3}{ \node[ln, minimum size=3.6mm, font=\tiny] (m\j) at (0.5*\j-1,1.05) {$\j$}; \node[fac, minimum width=4mm, minimum height=3.2mm] (g\j) at (0.5*\j-1,0.1) {}; }
        \draw[asg, line width=0.6pt] (m1) -- (g\a); \draw[asg, line width=0.6pt] (m2) -- (g\b); \draw[asg, line width=0.6pt] (m3) -- (g\c);
        \node[font=\scriptsize] at (0,-0.5) {$(\a\,\b\,\c)$};
      \end{scope}}
    \node[fdS, anchor=west] at (-3.9,-1.25) {each contributes $c_{a_1a_2a_3}$:\quad $i\cdot\bigl(-\frac{1}{3!}\bigr)\cdot3!\;c_{a_1a_2a_3} = -i\,c_{a_1a_2a_3}$};
  \end{scope}
\end{tikzpicture}
